%======================================================================
\section{Quantum channels that need nearly linear memory}\label{sec:channel}
%======================================================================

A finite set of relations becomes a quantum channel of fixed dimension: each relation contributes a few small unitary blocks, and the channel
evaluates the blocks in the regular representation of the group and traces out the group von Neumann algebra. In~\cite{DouglasMghirbiB} we showed
that the memory needed to apply such a channel repeatedly, releasing each output before the next input arrives, is bounded below by the size of
approximate representations of the relations. This section sharpens that link, so that it loses a single square root, and applies it to the
configuration lamps. The result is an explicit channel, in the closure of finite tracial channels, whose repeated use with logarithmic purity
needs memory $n^{1+o(1)}$: at least $n/(\log n)^{O(1)}$ for $n$ uses (Theorem~\ref{thm:main-channel}), and at most $O(n\log n)$, because the sofic
models of $\Gamma$ also give devices (Proposition~\ref{prop:sofic-upper}). More precisely, memory and purity trade at a fixed rate: a device
allowed purity $B$ needs memory about $n/B$, up to polylogarithmic factors, and exact devices attain this. At the optimal exchange rate the memory
is $n^{1/2+o(1)}$, attained by devices that keep one sofic model for all rounds (Theorem~\ref{thm:streaming}).

\subsection{Setting}\label{sec:channel-setting}

We recall the framework of~\cite{DouglasMghirbiB} and keep its notation.

\paragraph{Channels.} Channels act on $M_d=M_d(\mathbb C)$. The normalized Choi state of a channel $\Psi$ is
$\mathcal J(\Psi)=(\Psi\otimes\mathrm{id})(|\Omega\rangle\langle\Omega|)$, with $|\Omega\rangle=d^{-1/2}\sum_i|ii\rangle$, and
$\ddia(\Psi,\Psi')=\frac12\|\Psi-\Psi'\|_\diamond$ is the diamond distance. We write $\tau_M$ for the normalized trace on $M_M$. A channel $\Psi$ has a
\emph{finite tracial factorization} with bath dimension $M$ if $\Psi(X)=(\mathrm{id}\otimes\tau_M)[U(X\otimes I_M)U^*]$ for a unitary $U$ on
$\mathbb C^d\otimes\mathbb C^M$, and $\cK_d$ denotes the closure of the set of channels on $M_d$ with a finite tracial factorization of some bath
dimension. It is also the closure of the convex hull of that set~\cite[Proposition~3.8(ii)]{Musat2025}, and it consists of unital channels. Channels
with a finite tracial factorization are exactly the noisy operations of~\cite{HHO2003} from $M_d$ to $M_d$, and $\cK_d$ is their closure.

\paragraph{Causal services.} A \emph{causal service} of a channel $\Phi$ on $M_d$ for $n$ rounds consists of a memory register of $Q$ qubits in a fixed
initial state $\beta_0$ and, for each round, a fixed unitary on the input and the memory, after which the $d$-dimensional output is released; after
each round but the last, a prescribed set of memory qubits is exported for good and replaced by the same number of imported qubits in a fixed
state (possibly correlated within one step), independent of everything else. An observer chooses the inputs adaptively, possibly entangled with references of its own, and the
\emph{error} of the service is the largest half trace distance between the observer's final state in this experiment and in the ideal experiment
with $n$ independent uses of $\Phi$. The resources are the memory $Q$, the \emph{exchange} $J$ (the total number of exported qubits), and the
\emph{purity} $P+C$, where $P=Q-S(\beta_0)$ and $C$ is the sum over the imports of their number of qubits minus their entropy. Every retained
qubit counts; there is no free reset, measurement record or postselection.

\paragraph{Gadget channels.} Let $G$ be a group generated by a finite set $S$, and let $R$ be a finite set of words in $S^{\pm1}$, each of length at
least two and each representing $1$ in $G$. Write $r_i=\sigma_{i,1}\cdots\sigma_{i,\ell_i}$ ($1\le i\le r_0$) for the words of $R$, and put
$s=|S|$ and $\Lambda=\sum_i\ell_i$. The \emph{symbols} are the $2s$ letters $x^{\pm1}$, $x\in S$, and the proper prefixes
$p_{i,l}=\sigma_{i,1}\cdots\sigma_{i,l}$, $2\le l\le\ell_i-1$, one for each pair $(i,l)$; we also put $p_{i,1}=\sigma_{i,1}$ and let $p_{i,\ell_i}$
be the empty word. The \emph{triangles} are the triples $(p_{i,l-1},\sigma_{i,l},p_{i,l})$, $2\le l\le\ell_i$, and $(x,x^{-1},\varnothing)$,
$x\in S$. For unitaries $\lambda(w)$ assigned to the symbols, with $\lambda(\varnothing)=I$, let
\[
U_\lambda=\mathrm{diag}\bigl(I,(\lambda(w))_{w}\bigr)\ \oplus\ \bigoplus_{(x,y,z)}\frac1{\sqrt2}\begin{pmatrix}I&\lambda(y)\\ \lambda(x)&-\lambda(x)\lambda(y)\end{pmatrix},
\]
where $w$ runs over the symbols and $(x,y,z)$ over the triangles. Each triangle block equals $\mathrm{diag}(I,\lambda(x))\,(H\otimes I)\,\mathrm{diag}(I,\lambda(y))$
with $H$ the Hadamard matrix, so $U_\lambda$ is unitary for all unitary $\lambda$. There are $2s+\Lambda-2r_0$ symbols and $\Lambda-r_0+s$
triangles, so the system dimension is
\begin{equation}\label{eq:gadget-dimension}
d=1+(2s+\Lambda-2r_0)+2(\Lambda-r_0+s)=1+4s+3\Lambda-4r_0 .
\end{equation}
The \emph{gadget channel} $\Phi_G$ on $M_d$ evaluates each symbol $w$ in the left regular representation, $\lambda(w)=\lambda([w])$ with $[w]\in G$ the
element that $w$ represents, and traces out the group von Neumann algebra $L(G)$ with its canonical trace $\tau$:
$\Phi_G(X)=(\mathrm{id}\otimes\tau)[U_\lambda(X\otimes1)U_\lambda^*]$. Since the relations hold, the lower right entry of the block of a triangle
$(x,y,z)$ is $-\lambda([z])$.

We write the gadget as $U_\lambda=\sum_{a\in\mathcal T}c_aE_a\otimes\lambda([w_a])$, where $\mathcal T$ is the set of \emph{slots} (the nonzero
entries), $E_a$ is the matrix unit of slot $a$, $c_a\in\{\pm1,\pm2^{-1/2}\}$, and the word $w_a$ is empty, a symbol, or the concatenation $xy$ in the
lower right corner of a triangle. The elements $[w_a]$ are the \emph{labels}. With $A_g=\sum_{[w_a]=g}c_aE_a$, the matrices $A_g$ have disjoint
supports, $\Phi_G(X)=\sum_gA_gXA_g^*$ and $\mathcal J(\Phi_G)=\frac1d\sum_g|A_g\rangle\!\rangle\langle\!\langle A_g|$, because
$\tau(\lambda(g)\lambda(h)^*)=\delta_{gh}$. So the Choi rank $r_*$ of $\Phi_G$ is the number of distinct labels. Every label is carried by an
\emph{anchor}: a diagonal slot with coefficient $1$ whose row and column contain no other slot, namely the first diagonal entry for the label $1$
and a symbol entry otherwise, since the lower right entry of a triangle carries the label of a prefix symbol or $1$. Feeding one half of a maximally entangled
state on the $r_*$ anchor positions gives a pure probe whose output is flat of rank $r_*$, and $\log r_*$ is the largest entropy of a complementary
output of $\Phi_G$~\cite[Proposition~9.1]{DouglasMghirbiB}.

\paragraph{One round.} Let $U$ be a unitary on $\mathbb C^d\otimes\mathbb C^M$ and $\rho$ a state on $\mathbb C^M$, and put
\[
\Psi(X)=\mathrm{Tr}_M\,U(X\otimes\rho)U^*,\qquad T(\rho)=\mathrm{Tr}_d\,U(\tau_d\otimes\rho)U^* .
\]
Relative to a channel $\Phi$ with Choi support projection $\Pi_\Phi$, the \emph{support leakage} of $\Psi$ is
$\ell=\Tr[(I-\Pi_\Phi)\mathcal J(\Psi)]$. The \emph{entropy production} of the round is $\Delta S=S(T(\rho))-S(\rho)$, which is nonnegative because
$T$ is unital. A service produces such rounds: the next result, from~\cite{DouglasMghirbiB}, finds one in which both quantities are small.

\begin{theorem}[Sequential extraction~\cite{DouglasMghirbiB}]\label{thm:extraction}
Let a causal service of a channel $\Phi$ on $M_d$ for $n$ rounds have memory $Q$, error $\epsilon<1/2$ and $P+C\le B$, and put
\[
t_n=\frac{B+\epsilon+h_2(\epsilon)+\epsilon\log n}{n(1-\epsilon)} .
\]
Then some round $t$, with its unitary $U_t$ and the memory state $\rho_t$ entering it, defines a channel $\Psi_t(X)=\mathrm{Tr}_M\,U_t(X\otimes\rho_t)U_t^*$
with bath dimension $M=2^Q$, $\ddia(\Psi_t,\Phi)\le\epsilon/(1-\epsilon)$, and support leakage $\ell_t$ and entropy production $\Delta S_t$ satisfying
$\ell_t+\Delta S_t\le t_n$.
\end{theorem}

Since the bath has $Q$ qubits, $S(\rho_t)\le Q$. Every lower bound on memory below is a lower bound on the entropy of a memory state that realizes
one good round.

\subsection{The weighted one-round theorem}\label{sec:weighted-statement}

The following theorem is the engine of this section. It reads an approximate representation of the relations directly off the memory state of
one round, weighted by that state, and turns many cheaply commuting copies of $S_3$ into entropy of the memory. Transport words $u_1,\dots,u_N$
over $S^{\pm1}$ have \emph{length at most $L$} and \emph{pairwise areas at most $A$} if $|u_i|\le L$ for all $i$ and
\[
\Area_R\bigl([u_ixu_i^{-1},\,u_kyu_k^{-1}]\bigr)\le A\qquad(i\ne k,\ x,y\in\{a,b\}),
\]
where $\Area_R$ is the area in the free group on $S$, as in Section~\ref{sec:criterion}.

\begin{theorem}[Weighted one-round theorem]\label{thm:weighted}
Let $G$, $S$ and $R$ be as above, and suppose that $S$ contains $a$, $b$ and $j$, that $a^2$, $b^2$, $(ab)^3$ and $j^{-1}ab$ lie in $R$, and that
$j\ne1$ in $G$. Put $m=\max\{6,\ell_1,\dots,\ell_{r_0}\}$. Let $u_1,\dots,u_N$ be transport words of length at most $L$ and pairwise areas at most
$A$. Let $\Psi(X)=\mathrm{Tr}_M\,U(X\otimes\rho)U^*$ be one round with $\ddia(\Psi,\Phi_G)\le1/8$, support leakage $\ell$ relative to $\Phi_G$ and
entropy production $\Delta S$, and put $\nu=\ell+\Delta S$. If
\begin{equation}\label{eq:weighted-condition}
2^{25}\,m^2\sqrt d\,\sqrt N\,(A+L+1)\sqrt\nu\ \le\ 1,
\end{equation}
then $S(\rho)\ge N/16$.
\end{theorem}

Two features matter. First, the relator defect that the round reveals is of order $\sqrt\nu$, a single square root of the quantity that
Theorem~\ref{thm:extraction} controls, so a service with purity $O(\log n)$ is read at defect about $n^{-1/2}$. Second, there is no
tracialization: the argument works with the actual memory state $\rho$, which need not be flat. Replacing the round by a nearby finite tracial
mixture, as the general correspondence of~\cite{DouglasMghirbiB} does, costs a second square root, so that correspondence reads the profile at
defect about $n^{-1/4}$ instead of $n^{-1/2}$. The weighted route was
carried out in~\cite{DouglasMghirbiB} for a single gadget built from Houghton's group; Theorem~\ref{thm:weighted} holds for every relation
set that contains the relations of $S_3$ and a name for $ab$. It is proved in Section~\ref{sec:weighted-proof}.

\begin{remark}[The constant term]\label{rem:plus-one}
The term $+1$ in~\eqref{eq:weighted-condition} cannot be dropped. Take $N=1$, the empty transport word, so that $L=0$, and $A=0$, which is
allowed since there are no pairs. The gadget channel has a Stinespring dilation with a pure environment of dimension $r_*$, so there is a round
with $M=r_*$, $\rho$ pure and $\Psi=\Phi_G$ exactly. Then $\ddia(\Psi,\Phi_G)=0$ and the condition with $A+L$ in place of $A+L+1$ holds, since its left
side vanishes, but $S(\rho)=0<1/16$. With the $+1$, the condition forces $\nu\le2^{-50}m^{-4}d^{-1}$, while this round has $\ell=0$ and
$\Delta S=S(T(\rho))$, the entropy of the complementary output at the maximally mixed input. That output is diagonal with entries
$\|A_g\|_F^2/d\ge1/d$, and there are at least two labels, since $ab\ne1$; so $\Delta S\ge h_2(1/d)$, far above the threshold.
\end{remark}

Applied to a family of transport words that grows with a parameter, Theorem~\ref{thm:weighted} gives a lower bound on the entropy of one round
in terms of $\nu$ alone.

\begin{corollary}[Amplification]\label{cor:amplification}
Keep the assumptions of Theorem~\ref{thm:weighted} on $G$, $S$, $R$ and the round, and let $K\ge1$ and $\kappa\ge0$.
\begin{enumerate}[label=(\alph*),itemsep=1pt]
\item Suppose that for every integer $k\ge1$ there are $2^k$ transport words of length at most $L_k$ and pairwise areas at most $A_k$, with
$A_k+L_k+1\le Kk^\kappa$. Put $C_R=2^{50}m^4dK^2$. There is $y_0$ depending only on $\kappa$ such that
\[
S(\rho)\ \ge\ \frac{y}{32\,(\log y)^{2\kappa}}\qquad\text{for every }y\ge y_0\text{ with }C_R\,\nu\,y\le1 .
\]
\item Suppose instead that $A_k+L_k+1\le2^{K\sqrt k}$ for every $k\ge1$, and put $C'_R=2^{50}m^4d$. There is $y_1$ depending only on $K$ such that
\[
S(\rho)\ \ge\ \frac{y}{32}\,2^{-2K\sqrt{\log y}}\qquad\text{for every }y\ge y_1\text{ with }C'_R\,\nu\,y\le1 .
\]
\end{enumerate}
\end{corollary}

\begin{proof}
(a) Let $y_0\ge4$ be such that $k=\lfloor\log y-2\kappa\log\log y\rfloor\ge1$ for $y\ge y_0$. Then $k\le\log y$ and $2^k\le y(\log y)^{-2\kappa}$, so
$2^kk^{2\kappa}\le y$. With $N=2^k$, the left side of~\eqref{eq:weighted-condition} is at most
$2^{25}m^2\sqrt d\,K\,(2^kk^{2\kappa})^{1/2}\sqrt\nu\le(C_R\nu y)^{1/2}\le1$. So $S(\rho)\ge2^k/16$, and $2^k>y(\log y)^{-2\kappa}/2$.

(b) Let $y_1$ be such that $k=\lfloor\log y-2K\sqrt{\log y}\rfloor\ge1$ for $y\ge y_1$. Then $K\sqrt k\le K\sqrt{\log y}$, so
$2^k(A_k+L_k+1)^2\le2^{k+2K\sqrt{\log y}}\le y$, and~\eqref{eq:weighted-condition} holds with $N=2^k$ as in (a). Hence
$S(\rho)\ge2^k/16>y\,2^{-2K\sqrt{\log y}}/32$.
\end{proof}

\subsection{The channel of the configuration lamps}\label{sec:lamp-channel}

Theorem~\ref{thm:weighted} needs $a$, $b$ and a name for $ab$ among the generators. Let $R$ be the certificate of Table~\ref{tab:certificate}, in
the five generators $p,q,t,c,a$, with $b$ standing for the word $ac^4$. Let $S_\Phi=\{p,q,t,c,a,b,j\}$, where the new letters represent
$b=ac^4$ and $j=ab$ in $\Gamma$, and let $R_\Phi$ consist of the $25$ words of $R$ and the four words
\[
b^{-1}ac^4,\qquad b^2,\qquad (ab)^3,\qquad j^{-1}ab
\]
in the letters of $S_\Phi$. They represent $1$ in $\Gamma$, $a^2\in R$, and $j\ne1$ because $ab$ has order three (Lemma~\ref{lem:structure}). The
channel of the configuration lamps is the gadget channel $\Phi=\Phi_\Gamma$ of $R_\Phi$.

\paragraph{Dimension and rank.} Here $s=7$, $r_0=29$, and the total length is $\Lambda=303+6+2+6+3=320$, so~\eqref{eq:gadget-dimension} gives
\[
d=1+4\cdot7+3\cdot320-4\cdot29=873 ,
\]
with $276$ symbols and $298$ triangles; the longest word has length $24$, so $m=24$. The Choi rank $r_*$ is the number of distinct labels, at most
$1+276=277$. It is smaller: each of the $26$ words of $R_\Phi$ of length at least three has a longest proper prefix that represents the inverse of
its last letter, so these $26$ prefix symbols carry labels already carried by letters, and $a^2=b^2=t^2=1$ identifies the labels of
$a^{-1},b^{-1},t^{-1}$ with those of $a,b,t$. Hence $2\le r_*\le248$, where the lower bound holds because $1$ and $j$ are distinct labels.

\paragraph{Transport.} For $N\ge1$, the proof of Theorem~\ref{thm:main-group} provides $2^N$ words $T_u$, one for each configuration $u$
supported in the first $N$ points of ray~$1$, with $|T_u|\le21N\max(3,N)\le63N^2$ (Theorem~\ref{thm:pair-area}) and
$\Area_{R_\Phi}([T_uxT_u^{-1},T_wyT_w^{-1}])\le\Lambda_N+4\le C_1N^{\kappa}$ for $u\ne w$ and $x,y\in\{a,b\}$, where
\[
p_H=\log107+4,\qquad \kappa=2+p_H<12.75,
\]
and $C_1$ is absolute; the letter $b$ is converted to $ac^4$ by the word $b^{-1}ac^4$, at most four times, and areas with respect to
$R_\Phi\supseteq R$ are no larger than with respect to $R$. So the hypothesis of Corollary~\ref{cor:amplification}(a) holds with this $\kappa$ and
$K=C_1+64$.

\begin{theorem}[Theorem B]\label{thm:main-channel}
The channel $\Phi$ on $M_{873}$ lies in $\cK_{873}$, and its Choi rank satisfies $2\le r_*\le248$. Let $\kappa=2+p_H=6+\log107$. There are
constants $c,c_0,K>0$ and $n_0$ such that the following hold; all services have error at most $1/16$.
\begin{enumerate}[label=(\alph*),itemsep=1pt]
\item \emph{Memory against purity.} For $n\ge n_0$ and $\log n\le B\le c_0n$, every causal service of $\Phi$ for $n$ rounds with purity
$P+C\le B$ has, at any exchange,
\[
Q\ \ge\ c\,\frac{n}{B\,(\log n)^{2\kappa}} .
\]
For every $n\ge1$ and $0<B\le n$ there is an exact causal service, with error $0$, purity $P+C<B$, exchange $J=(n-1)Q$ and memory
$Q\le K(n/B)\log(2n/B)$. In particular, for every $A_0\ge0$ there is $c'>0$ such that services with $P+C\le A_0\log n$ have
$Q\ge c'n/(\log n)^{2\kappa+1}$ for all large $n$, where $2\kappa+1=2p_H+5<26.49$, and there are services with $P=C=0$ and $Q\le Kn\log n$.
\item \emph{The rank rate.} For $n\ge n_0$, every causal service of $\Phi$ for $n$ rounds at the rank rate, $J\le\frac n2\log r_*$, has, whatever
its purity,
\[
Q\ \ge\ c\,\frac{\sqrt n}{(\log n)^{\kappa}} .
\]
For every $n\ge2$ there is a causal service at the rank rate with $C=0$ and $P,Q\le K\sqrt{n\log n}$.
\end{enumerate}
So memory and purity trade at a fixed rate: for $n\ge n_0$ and $\log n\le B\le c_0n$ the least memory of a service with purity at most $B$ lies
between $c\,n/(B(\log n)^{2\kappa})$ and $K(n/B)\log(2n/B)$, and the product of memory and purity is $n^{1+o(1)}$. At the rank rate the least
memory is $n^{1/2+o(1)}$, and it is attained with purity of the same order: the rank rate sits where memory and purity are equal.
\end{theorem}

The exponent $2\kappa$ in (a) is the unitary exponent of Theorem~\ref{thm:main-group}(c): a service with purity $B$ is read at relator defect about
$(B/n)^{1/2}$. The lower bound in (b) needs no assumption on purity, because at the rank rate purity is paid for with memory
(Proposition~\ref{prop:purity-budget}): there $P+C<3Q+6$, and the argument of (a) with $B$ of order $Q$ gives $Q^2\gtrsim n$ up to logarithms. The
exact services in (a) refresh the whole memory in every round, so their exchange is far above the rank rate; the services in (b) keep a sofic
model in memory for all $n$ rounds and pay for its bad points with purity (Theorem~\ref{thm:streaming}). We prove Theorem~\ref{thm:main-channel} in Section~\ref{sec:main-channel-proof}, after the purity budget
and the upper bound.

\begin{remark}[The grid host]\label{rem:grid-channel}
The grid host $\Gamma_2$ of Remark~\ref{rem:grid} should give a stronger lower bound with the same upper bound; both rest on sketches. Its certificate has $7$ generators and $77$
relations of total length $5445$ and maximum length $224$. Adding $b$ and $j$ and the same four words gives $s=9$, $r_0=81$ and $\Lambda=5462$, so
its gadget channel $\Phi_2$ acts on $M_d$ with
\[
d=1+4\cdot9+3\cdot5462-4\cdot81=16099,
\]
has $m=224$, and has Choi rank at most $1+18+5300-78-3=5238$ by the count above ($78$ of its $81$ words have length at least three). With the
transport of Remark~\ref{rem:grid}, words of length $O(N^{3/2})$ and pair areas $O(N^2(\delta_H(C\sqrt N)+N))$, the hypothesis of
Corollary~\ref{cor:amplification}(a) holds with $\kappa=2+p_H/2$, and part (b) of that corollary holds with Lee's exponential Dehn
bound~\cite{Lee2012} alone. The proof below then gives, for $\Phi_2$: memory at least $c\,n/(\log n)^{p_H+5}$ at purity $O(\log n)$, where
$p_H+5<15.75$; memory at least $n\,2^{-C\sqrt{\log n}}$ at purity $O(\log n)$ using only Lee's bound; and memory at least
$c\sqrt n/(\log n)^{2+p_H/2}$ at the rank rate, whatever the purity. These rest on the pair-area bound for $\Gamma_2$, whose proof
Remark~\ref{rem:grid} only sketches. By the two-stage hashing sketched in Remark~\ref{rem:grid}, $\log\sigma(r)=O(r\log r)$ for $\Gamma_2$ as well, and so an upper bound of order
$n\log n$. For $\Gamma$ itself, Lee's bound alone gives pair areas $2^{O(N)}$, and the argument gives only memory $n^{c}$ for some $c>0$.
\end{remark}

\subsection{A purity budget at the rank rate}\label{sec:purity-budget}

At the rank rate a device cannot buy memory with purity: the purity it can consume is linear in its memory. We sharpen the budget $14Q+3$
of~\cite{DouglasMghirbiB} to $3Q+6$ by combining its rank count with a one-sided second-moment bound.

\begin{proposition}[Purity budget]\label{prop:purity-budget}
Let $\Phi$ be a channel on $M_d$, and suppose that some pure state $\psi_*$ of the input and a reference has output
$\omega=(\Phi\otimes\mathrm{id})(\psi_*)$ flat of rank $r$. Every causal service of $\Phi$ for $n$ rounds with error $\epsilon\le1/16$ and exchange
$J\le\frac n2\log r$ satisfies
\[
P+C\ <\ 3Q+6 .
\]
\end{proposition}

Every gadget channel has such a probe with $r=r_*$ (Section~\ref{sec:channel-setting}). The proof uses the variance of the surprisal of a
state with small deficit, which is~\cite[Lemma~5.9]{DouglasMghirbiB}; we include the short argument.

\begin{lemma}[{Surprisal variance~\cite[Lemma~5.9]{DouglasMghirbiB}}]\label{lem:surprisal}
Let $\beta$ be a state on $q\ge1$ qubits with spectrum $(\beta_i)$ and deficit $D=q-S(\beta)$. The surprisal $-\log\beta_i$, drawn with probability
$\beta_i$, has variance at most $(q+2/\ln2)\,D$.
\end{lemma}

\begin{proof}
Put $\varphi(x)=x\ln x-x+1\ge0$. We claim $x(\ln x)^2\le(2+\ln_+x)\varphi(x)$ for $x\ge0$. For $x=e^t$ with $t\ge0$ the difference of the two sides
is $f(t)=e^t(t-2)+t+2$, with $f(0)=f'(0)=0$ and $f''(t)=te^t\ge0$. For $t\le0$ it is $g(t)=2(te^t-e^t+1)-t^2e^t$, with $g(0)=0$ and
$g'(t)=-t^2e^t\le0$; and $x=0$ is clear. Now let $Q_0=2^q$ and $x_i=Q_0\beta_i\le Q_0$, the index running over all $Q_0$ eigenvalues. Then
$\frac1{Q_0}\sum_ix_i=1$ and $D\ln2=\sum_i\beta_i\ln(Q_0\beta_i)=\frac1{Q_0}\sum_i\varphi(x_i)$. The variance is at most
\[
\sum_i\beta_i\bigl(\log(Q_0\beta_i)\bigr)^2=\frac1{Q_0(\ln2)^2}\sum_ix_i(\ln x_i)^2\le\frac{2+q\ln2}{(\ln2)^2}\,D\ln2=\Bigl(q+\frac2{\ln2}\Bigr)D .
\]
\end{proof}

\begin{proof}[Proof of Proposition~\ref{prop:purity-budget}]
Put $b=P+C$ and suppose, for a contradiction, that $b\ge3Q+6$.

\emph{Preloading.} The imports are fixed states, independent of everything else, and the exported qubits are never used again. So the observer's
final state does not change if every import is prepared at the start in a register of its own, and each replacement swaps it with the exported
qubits, which then stay in the device. The device becomes a register of $L=Q+J$ qubits in the product state
$\sigma_0=\beta_0\otimes\bigotimes_f\omega_f$ of the initial memory state and the imported states, and it applies a fixed unitary in each round.

\emph{The surprisal.} Let $(\sigma_i)$ be the spectrum of $\sigma_0$ and let $Z=L+\log\sigma_i$, with $i$ drawn with probability $\sigma_i$. Then
$\mathbb EZ=L-S(\sigma_0)=b$. Also $Z$ is the sum of the independent variables $Z_f=q_f+\log\beta^{(f)}_i$ attached to the factors of $\sigma_0$,
where $q_f$ is the number of qubits of the factor and $\beta^{(f)}$ its spectrum. Every factor has at most $Q$ qubits, since an import replaces
exported memory qubits, and the deficits $b_f=\mathbb EZ_f$ of the factors add up to $b$. By Lemma~\ref{lem:surprisal},
\[
\mathrm{Var}\,Z=\sum_f\mathrm{Var}\,Z_f\le\Bigl(Q+\frac2{\ln2}\Bigr)b .
\]

\emph{A tail bound from the rank.} Let the observer feed the input halves of $n$ copies of $\psi_*$ and keep the references. Its final state is
$\tau_O=\sum_i\sigma_i\tau_i$, where $\tau_i$ is its final state when the register starts in the eigenvector of $\sigma_i$. Each $\tau_i$ is a
marginal of a pure state of the observer's system and the $L$-qubit register, so $\rank\tau_i\le2^L$. Fix a real number $k$ and let $I_k$ be the
set of indices with $\sigma_i\ge2^{k-L}$, so that $|I_k|\le2^{L-k}$. Let $E$ be the support projection of $\sum_{i\in I_k}\tau_i$. Then
$\rank E\le2^{2L-k}$ and $\Tr(E\tau_O)\ge\sum_{i\in I_k}\sigma_i=\Pr(Z\ge k)$. The ideal final state $\omega^{\otimes n}$ is flat of rank $r^n$, so
$\Tr(E\omega^{\otimes n})\le2^{2L-k}r^{-n}\le2^{2Q-k}$, using $2J\le n\log r$. Since the error is at most $\epsilon$,
\begin{equation}\label{eq:rank-tail}
\Pr(Z\ge k)\ \le\ \epsilon+2^{2Q-k}\qquad(k\in\mathbb R).
\end{equation}

\emph{The contradiction.} Put $u=b-2Q-3$, so that $u\ge b/3+1>0$. By the one-sided Chebyshev inequality,
$\Pr(Z\ge b-u)\ge u^2/(\mathrm{Var}\,Z+u^2)$, while~\eqref{eq:rank-tail} with $k=b-u=2Q+3$ gives $\Pr(Z\ge b-u)\le\frac1{16}+\frac18=\frac3{16}$. Hence
$13u^2\le3\,\mathrm{Var}\,Z$, and with $u>b/3$,
\[
\frac{13}9\,b^2<3\Bigl(Q+\frac2{\ln2}\Bigr)b,\qquad\text{that is}\qquad b<\frac{27}{13}\Bigl(Q+\frac2{\ln2}\Bigr)<2.08\,Q+5.993 ,
\]
which contradicts $b\ge3Q+6$.
\end{proof}

\subsection{An upper bound from sofic models}\label{sec:sofic-upper}

Sofic models of the labels give devices. The point is that the triangle blocks can be made exactly unitary for arbitrary permutations, and that at
most points of a model the pure bath vector reproduces the target channel exactly, not approximately. Recall from Section~\ref{sec:criterion}
that $\sigma_E(r)$ is the least $M$ for which some $f:E\to\Sym(\Omega)$, $|\Omega|=M$, with $f(1)=\mathrm{id}$, is a
\emph{$(1-r^{-1})$-expansive $r^{-1}$-morphism}: $f(x)f(y)$ and $f(xy)$ differ on at most $M/r$ points whenever $x,y,xy\in E$, and $f(x)$ and
$f(y)$ agree on at most $M/r$ points whenever $x\ne y$ in $E$.

\begin{proposition}[Sofic models give devices]\label{prop:sofic-upper}
Let $G$, $S$ and $R$ be as in Section~\ref{sec:channel-setting}, with $t_\triangle$ triangles and $r_*$ labels, let $E$ be a chunk of $G$ containing
all labels, and put $K_G=t_\triangle+\binom{r_*}2$. Every $(1-r^{-1})$-expansive $r^{-1}$-morphism $f:E\to\Sym(\Omega)$, $|\Omega|=M$, gives a unitary $U_f$ on
$\mathbb C^d\otimes\mathbb C^\Omega$ whose tracial channel $\Phi_f(X)=(\mathrm{id}\otimes\tau_M)[U_f(X\otimes I)U_f^*]$ satisfies
\[
\Phi_f=(1-\vartheta)\Phi_G+\vartheta\,\Xi
\]
for some channel $\Xi$ and some $0\le\vartheta\le K_G/r$. Consequently, for all $n\ge1$ and $0<\epsilon\le1$ with $\sigma_E(2K_Gn/\epsilon)<\infty$,
$\Phi_G$ has
\begin{enumerate}[label=(\alph*),itemsep=1pt]
\item a causal service with error at most $\epsilon$, $P=C=0$, $J=(n-1)Q$ and $Q=\bigl\lceil\log\sigma_E(2K_Gn/\epsilon)+\log(2n/\epsilon)\bigr\rceil$;
\item a causal service with error at most $\epsilon$, $J=P=C=0$ and $Q=\bigl\lceil n\log\sigma_E(2K_Gn/\epsilon)+\log(2/\epsilon)\bigr\rceil$;
\end{enumerate}
and, for all $n\ge1$ and $0<B\le n$ with $\sigma_E(4K_Gn/B)<\infty$,
\begin{enumerate}[label=(\alph*),itemsep=1pt,resume]
\item an exact causal service, with error $0$, purity $P+C<B$, $J=(n-1)Q$ and $Q=\bigl\lceil\log\sigma_E(4K_Gn/B)+\log(4n/B)\bigr\rceil$.
\end{enumerate}
\end{proposition}

\begin{proof}
\emph{The unitary.} For a symbol $w$ let $P_w$ be the permutation matrix of $f([w])$ on $\mathbb C^\Omega$, so that symbols with the same label get the
same matrix, and let $U_f$ be the gadget $U_\lambda$ with $\lambda(w)=P_w$. Each triangle block is
$\frac1{\sqrt2}\bigl(\begin{smallmatrix}I&P_y\\P_x&-P_xP_y\end{smallmatrix}\bigr)=\mathrm{diag}(I,P_x)(H\otimes I)\mathrm{diag}(I,P_y)$, so $U_f$ is unitary
although no relation holds among the $P_w$. Writing $\pi_a$ for the permutation of $\Omega$ in slot $a$, namely $f([w_a])$ for every slot except the
lower right corner of a triangle $(x,y,z)$, where it is $f([x])f([y])$, we have $U_f=\sum_ac_aE_a\otimes P_{\pi_a}$.

\emph{Pure bath points.} Since $\tau_M$ is the average of the vector states of the points of $\Omega$, $\Phi_f=\frac1M\sum_{\omega\in\Omega}\Phi_\omega$ with
$\Phi_\omega(X)=\mathrm{Tr}_\Omega\,U_f(X\otimes|\omega\rangle\langle\omega|)U_f^*$. The isometry $|\psi\rangle\mapsto U_f(|\psi\rangle\otimes|\omega\rangle)$ is
$\sum_ac_aE_a\otimes|\pi_a(\omega)\rangle$, so
\begin{gather*}
\Phi_\omega(X)=\sum_{a,b\in\mathcal T}c_a\overline{c_b}\;\mathbf 1[\pi_a(\omega)=\pi_b(\omega)]\;E_aXE_b^*,\\
\Phi_G(X)=\sum_{a,b\in\mathcal T}c_a\overline{c_b}\;\mathbf 1\bigl[[w_a]=[w_b]\bigr]\;E_aXE_b^*,
\end{gather*}
the second because $\tau(\lambda(g)\lambda(h)^*)=\delta_{gh}$. Call $\omega$ \emph{good} if $f([x])f([y])(\omega)=f([z])(\omega)$ for every triangle
$(x,y,z)$ and $f(g)(\omega)\ne f(h)(\omega)$ for all distinct labels $g,h$. At a good point $\pi_a(\omega)=f([w_a])(\omega)$ for every slot, since
$[w_a]=[z]$ at the lower right corner of a triangle $(x,y,z)$; so $\pi_a(\omega)=\pi_b(\omega)$ exactly when $[w_a]=[w_b]$, and $\Phi_\omega=\Phi_G$.

\emph{Bad points.} The labels of a triangle satisfy $[x][y]=[z]$ with all three in $E$, so $f([x])f([y])$ and $f([z])$ differ on at most $M/r$ points;
and distinct labels $g,h\in E$ have $f(g)(\omega)=f(h)(\omega)$ for at most $M/r$ points. So the fraction $\vartheta$ of bad points is at most $K_G/r$,
and $\Phi_f=(1-\vartheta)\Phi_G+\vartheta\,\Xi$, where $\Xi$ is the average of $\Phi_\omega$ over the bad points (any channel if there are none).

\emph{Exchange.} Let $r=2K_Gn/\epsilon$, let $f$ have degree $M=\sigma_E(r)$, and let $Q$ be as in (a), so that $M\le2^Q\epsilon/(2n)$. Split the
$Q$-qubit memory as $(\mathbb C^\Omega\otimes\mathbb C^{k_0})\oplus\mathbb C^{g_0}$ with $k_0=\lfloor2^Q/M\rfloor$ and $g_0=2^Q-k_0M<M$. In every round
apply $U_f\otimes I_{k_0}$ on the first summand and the identity on the second, and after every round but the last export all $Q$ memory qubits and
import $Q$ fresh qubits in the maximally mixed state; the initial state is maximally mixed as well, so $P=C=0$. Every round then meets a fresh
flat memory, and the observer faces $n$ independent uses of the channel $(1-g_0/2^Q)\Phi_f+(g_0/2^Q)\mathrm{id}$, which is within
$\vartheta+g_0/2^Q\le\epsilon/(2n)+\epsilon/(2n)$ of $\Phi_G$ in diamond distance, by convexity. A hybrid argument over the rounds bounds the error by $\epsilon$.

\emph{Exact service.} Let $0<B\le n$, put $\eta=B/(4n)\le\frac14$ and $r=K_G/\eta$, let $f$ have degree $M=\sigma_E(r)$, and let $Q$ be as in (c), so
that $M\le\eta2^Q$. Let $\Omega_0\subseteq\Omega$ be the set of good points, so that $|\Omega_0|\ge(1-\eta)M$. Split the memory as
$(\mathbb C^\Omega\otimes\mathbb C^{k_0})\oplus\mathbb C^{g_0}$ as in the exchange case, with $k_0=\lfloor2^Q/M\rfloor$ and $g_0<M\le\eta2^Q$, and let
$\beta$ be the maximally mixed state on $\mathrm{span}\{|\omega\rangle:\omega\in\Omega_0\}\otimes\mathbb C^{k_0}$. Since $k_0M=2^Q-g_0$, its deficit is
\[
Q-\log(|\Omega_0|k_0)=\log\frac{2^Q}{2^Q-g_0}+\log\frac M{|\Omega_0|}\le2\log\frac1{1-\eta}\le\frac{8\eta}{3\ln2}<\frac Bn .
\]
Start the memory in $\beta$, apply $U_f\otimes I_{k_0}$ on the first summand in every round, and after every round but the last export all $Q$
memory qubits and import $Q$ fresh qubits in the state $\beta$. The $n$ registers used have total deficit $P+C<B$. In every round the memory is a
fresh copy of $\beta$, a uniform mixture of vector states $|\omega\rangle\otimes|j\rangle$ with $\omega$ good, and $\Phi_\omega=\Phi_G$ at good points;
so every round applies exactly $\Phi_G$, independently of the other rounds, and the error is $0$.

\emph{No exchange.} Let $f$ be as in the exchange case and $Q$ as in (b), so that $M^n\le2^Q\epsilon/2$. Split the memory as
$(\mathbb C^{\Omega^n}\otimes\mathbb C^{k_0})\oplus\mathbb C^{g_0}$ with $k_0=\lfloor2^Q/M^n\rfloor$ and $g_0<M^n$, start it maximally mixed, and in round $s$
apply $U_f$ to the input and the $s$-th factor $\mathbb C^\Omega$ of the first summand, and the identity on the second. With probability
$1-g_0/2^Q\ge1-\epsilon/2$ the memory starts in the first summand, uniformly distributed, and then the $n$ factors are independent flat baths, each
used once: the observer faces $n$ independent uses of $\Phi_f$, which is within $n\vartheta\le\epsilon/2$ of the ideal experiment. The remaining
branch has weight at most $\epsilon/2$, so the error is at most $\epsilon$.
\end{proof}

In particular, if $\sigma_E(r)$ is finite for every $r$, as it is when $G$ is sofic, then $\Phi_G\in\cK_d$, since
$\ddia(\Phi_f,\Phi_G)\le K_G/r$ for models at every accuracy. This is a quantitative form, for sofic groups, of the fact that gadget channels of
hyperlinear groups lie in $\cK_d$~\cite{DouglasMghirbiB}.

\subsection{Streaming at the rank rate}\label{sec:streaming}

The services of Proposition~\ref{prop:sofic-upper} refresh the whole memory in every round. At the rank rate the memory has to be kept, and a sofic
model cannot simply be reused, since its bad points would then accumulate error. The following theorem keeps the model for all rounds and pays
for its bad points in purity instead. It uses the good points of a model as an exact \emph{rectangular} dilation, whose input bath is smaller than
its output bath, and lets the part of the memory in use grow slightly in every round. The mixing and export steps are those of the rolling
reservoir of~\cite[Appendix~A]{DouglasMghirbiB}.

\begin{theorem}[Streaming at the rank rate]\label{thm:streaming}
Let $\Phi$ be a channel on $M_d$ with linearly independent Kraus operators $A_1,\dots,A_r$, $r\ge2$, and let $1\le k\le k'$. Suppose that
$W=\sum_iA_i\otimes B_i$, with $B_i:\mathbb C^k\to\mathbb C^{k'}$, is an isometry from $\mathbb C^d\otimes\mathbb C^k$ to $\mathbb C^d\otimes\mathbb C^{k'}$ and that
\begin{equation}\label{eq:rectangular}
\Tr(B_i^*B_j)=k\,\delta_{ij},\qquad \sum_iB_i^*B_i=r\,I_k .
\end{equation}
Then for all $n\ge2$ and $0<\epsilon<1$ there is a causal service of $\Phi$ for $n$ rounds with error at most $\epsilon$, $C=0$,
$J=\lfloor\frac{n-1}2\log r\rfloor$ and
\[
P\le n\log\frac{k'}k+2,\qquad Q\le\log k+n\log\frac{k'}k+O_r\Bigl(\log\frac n\epsilon+\log(2+\log k)\Bigr).
\]
\end{theorem}

The first identity in~\eqref{eq:rectangular} says that $W$ with a flat input bath realizes $\Phi$ exactly:
$\Phi(X)=\mathrm{Tr}_{k'}[W(X\otimes I_k/k)W^*]=\sum_iA_iXA_i^*$. The second says that every Kraus history costs the same. For a channel with a
flat probe of rank $r$ it follows from $W^*W=I$ alone: apply $\Tr(\rho_*\,\cdot\,)$ to the first factor of $\sum_{i,j}A_i^*A_j\otimes B_i^*B_j=I$,
where $\rho_*$ is the input marginal of the probe, for which $\Tr(\rho_*A_i^*A_j)=\delta_{ij}/r$.

\begin{proof}
We use the frames of~\cite[Appendix~A]{DouglasMghirbiB}. A \emph{frame} is a linear map $F:\mathbb C^M\to\mathbb C^D\otimes\mathbb C^S$; its columns
are the vectors, in the memory $\mathbb C^D$ and an inaccessible exterior $\mathbb C^S$, attached to the Kraus histories. The exterior holds the
purifier of the initial memory, the exported qubits and the purifiers of the imported ones. Put $g(F)=\|F^*F\|$ and
$\alpha(F)=\|\mathrm{Tr}_{\mathbb C^D}FF^*\|$. Two facts are used.
\begin{itemize}[leftmargin=1.2em,itemsep=1pt]
\item \emph{Concentration}~\cite[Lemma~A.1]{DouglasMghirbiB}. For Hermitian $X$ on $\mathbb C^D$ and Haar-random $V\in U(D)$, the matrix
$Z=F^*((V^*XV)\otimes I)F-\tau_D(X)F^*F$ satisfies $\Pr\{\|Z\|\ge s\}\le2M\exp[-Ds^2/(4\alpha(F)g(F)\|X\|^2)]$.
\item \emph{Export}~\cite[Lemma~A.4]{DouglasMghirbiB}. If $\mathbb C^D=\mathbb C^{D/e}\otimes\mathbb C^e$, there is a unitary of $\mathbb C^D$ after
which exporting the factor $\mathbb C^e$ and importing a maximally mixed $e$-dimensional register give a frame $F'$ with $F'^*F'=F^*F$ and, for every
$\eta>0$, $\alpha(F')\le(1+\eta)\alpha(F)/e^2+4e^2g(F)\ln(8Se^2)/(\eta D)$.
\end{itemize}

\emph{Schedule and dimensions.} Put $h=\log r$, $E_p=2^{\lceil h/2\rceil}$, $J_t=\lfloor th/2\rfloor$, $j_t=J_t-J_{t-1}\le\lceil h/2\rceil$ and
$e_t=2^{j_t}$, a divisor of $E_p$. Choose $D_0$, a multiple of $kE_p$, with $D_0\ge nkE_p$ and
\begin{equation}\label{eq:D0}
D_0^2\ \ge\ \frac{32\,A_*\,r^4n^2\ln(8r^{n+2})}{\epsilon^2},\qquad\text{where}\quad A_*=12\bigl(1+8E_p^2n^2L_*\bigr),\quad L_*=\ln(8D_0r^nE_p^2).
\end{equation}
Put $q_0=k'/k$. For $t\ge1$ let $\widetilde D_t=q_0D_{t-1}$, a multiple of $k'E_p$, and let $D_t$ be the least multiple of $kE_p$ with
$D_t\ge\widetilde D_t$. Then $D_t\le q_0D_{t-1}+kE_p$, so $D_n\le q_0^n(D_0+nkE_p)\le2q_0^nD_0$. The memory has $Q=\lceil\log D_n\rceil$ qubits, of which the last $\lceil h/2\rceil$
form a \emph{port}. A space whose dimension $D'$ is a multiple of $E_p$ is realized as a subspace of the first $Q-\lceil h/2\rceil$ qubits of dimension
$D'/E_p$, tensored with the port; the \emph{active space} of round $t$ is such a space of dimension $D_{t-1}$, identified with
$\mathbb C^k\otimes\mathbb C^{D_{t-1}/k}$.

\emph{The device.} The memory starts maximally mixed on the active space of round $1$, of dimension $D_0$. In round $t$ the device applies a unitary
$V_t$ of the active space and then $W\otimes I_{D_{t-1}/k}$, which maps the input and the active space isometrically onto the output and a space of
dimension $\widetilde D_t$; it releases the output, then applies a unitary $V'_t$ of that space (it acts on the memory only, so it could equally act before the release), and, if $t<n$, exports $j_t$ port qubits and imports $j_t$
qubits in the maximally mixed state in their place. The resulting space of dimension $\widetilde D_t$ lies in the active space of round $t+1$. These
maps are completed arbitrarily to unitaries of the input and the whole memory; the completion never acts, since the memory never leaves the
active spaces. So $C=0$, $J=J_{n-1}$, and $P=Q-\log D_0\le n\log q_0+2$.

\emph{Frames.} Purify the initial memory: its frame $F_0$ is one unit vector, with $\alpha(F_0)=1/D_0$. Let $F$ be the frame before round $t$, with one
column $F_y$ for each history $y\in[r]^{t-1}$. After the call the frame has the columns $(B_i\otimes I)(V_t\otimes I)F_y$, indexed by the histories
$(y,i)$. Its Gram matrix has the blocks $F^*\bigl((V_t^*(B_i^*B_j\otimes I)V_t)\otimes I\bigr)F$, whose Haar means are
$\tau(B_i^*B_j\otimes I)F^*F=\delta_{ij}F^*F$ by the first identity in~\eqref{eq:rectangular}. By the second identity and cyclicity of the partial trace
over the memory, the exterior marginal after the call is exactly $r$ times the one before, whatever $V_t$ is; so the call multiplies $\alpha$ by $r$.

\emph{Choice of the unitaries.} We choose $V_1,V'_1,V_2,\dots$ in turn so that the frame $F_t$ after round $t$ satisfies
\begin{equation}\label{eq:streaming-induction}
\|F_t^*F_t-I\|\le t\epsilon/n,\qquad \alpha(F_t)\le A_*/D_0 .
\end{equation}
This holds for $t=0$. Suppose it holds for $t-1$; then $g(F_{t-1})\le2$. Since $B_i^*B_i\le rI$, the Hermitian and anti-Hermitian parts of the $B_i^*B_j$
have norm at most $r$. Apply the concentration lemma to these $2r^2$ operators, tensored with $I_{D_{t-1}/k}$, with $D=D_{t-1}\ge D_0$, at most $r^n$
columns and level $s=\epsilon/(2rn)$: by~\eqref{eq:D0} the failure probabilities add up to at most
$4r^{n+2}\exp[-D_0^2\epsilon^2/(32A_*r^4n^2)]\le\frac12$. So some $V_t$ makes every block of the new Gram matrix deviate from its mean by at most
$\epsilon/(rn)$, hence the $r\times r$ block matrix by at most $\epsilon/n$, and the Gram matrix after the call is within $t\epsilon/n$ of $I$; exports and
imports preserve it. Choose $V'_t$ by the export lemma with $e=e_t$, $\eta=1/n$ and $D=\widetilde D_t\ge D_0$. The exterior has dimension
$S\le D_0\,2^{2J_{t-1}}\le D_0r^n$, so $\ln(8Se_t^2)\le L_*$, and
\[
\alpha(F_t)\ \le\ \Bigl(1+\frac1n\Bigr)\frac{r\,\alpha(F_{t-1})}{e_t^2}+\frac{8E_p^2nL_*}{D_0}.
\]
Put $c_t=2^{2J_t}r^{-t}$, so that $c_0=1$, $\frac14<c_t\le1$ and $c_tr/e_t^2=c_{t-1}$. Multiplying by $c_t$ and unrolling gives
$c_t\alpha(F_t)\le(1+\frac1n)^t(1+8E_p^2n^2L_*)/D_0$, hence $\alpha(F_t)\le4e(1+8E_p^2n^2L_*)/D_0\le A_*/D_0$. In round $n$ there is no export, and only the
Gram bound is needed.

\emph{Error.} Purify the observer. In the ideal experiment, with the Stinespring isometry $\xi\mapsto\sum_iA_i\xi\otimes|i\rangle$ in each round, its
final vector is $\sum_y\psi_y\otimes|y\rangle$; in the service it is $\sum_y\psi_y\otimes F_n|y\rangle$ with the same vectors $\psi_y$, because the device
acts on the input and the memory through $W=\sum_iA_i\otimes B_i$ and unitaries of the memory, and the observer never touches the memory or the
exterior. The polar part $T=F_n(F_n^*F_n)^{-1/2}$ is an isometry with $\|F_n-T\|=\|(F_n^*F_n)^{1/2}-I\|\le\epsilon$. So the actual final vector is within
$\epsilon$ of the ideal one mapped by $I\otimes T$, and after tracing out the memory and the exterior the observer's final states are within
$\epsilon$ in trace distance.

\emph{Memory.} Finally $Q\le\log D_n+1\le\log D_0+n\log q_0+2$. The conditions on $D_0$ hold for $D_0=kE_p\lceil K_r(n/\epsilon)^4(1+\ln k)\rceil$ with $K_r$
large enough in terms of $r$, and then $\log D_0=\log k+O_r(\log(n/\epsilon)+\log(2+\log k))$.
\end{proof}

\begin{corollary}[The lamp channel at the rank rate]\label{cor:streaming-lamps}
There is $K$ such that for every $n\ge2$ the channel $\Phi$ of Theorem~\ref{thm:main-channel} has a causal service for $n$ rounds with error at most
$1/16$, $C=0$, $J=\lfloor\frac{n-1}2\log r_*\rfloor$ and $P,Q\le K\sqrt{n\log n}$.
\end{corollary}

\begin{proof}
Let $E$ be the chunk of labels of Proposition~\ref{prop:sofic-upper}, let $s\ge2K_G$, and let
$f:E\to\Sym(\Omega)$ be a $(1-s^{-1})$-expansive $s^{-1}$-morphism of degree $M=\sigma_E(s)$. Let $\Omega_0$ be its set of good points, so that
$k=|\Omega_0|\ge(1-K_G/s)M$, and for each label $g$ let $B_g:\mathbb C^{\Omega_0}\to\mathbb C^\Omega$ send $|\omega\rangle$ to $|f(g)\omega\rangle$. At a good point
$\omega$ every slot of $U_f$ acts through the permutation of its label, so $U_f(\xi\otimes|\omega\rangle)=\sum_gA_g\xi\otimes|f(g)\omega\rangle$, and
$W=\sum_gA_g\otimes B_g$ is the restriction of the unitary $U_f$ to $\mathbb C^d\otimes\mathbb C^{\Omega_0}$, an isometry. Each $B_g^*B_g$ is $I_k$, and
$\Tr(B_g^*B_h)=0$ for $g\ne h$ because distinct labels do not collide at good points; the $A_g$ are linearly independent, having disjoint
supports. So Theorem~\ref{thm:streaming} applies with $r=r_*$ and $k'=M$. Here $\log(M/k)\le-\log(1-K_G/s)\le2K_G/(s\ln2)$, and
$\log k\le\log\sigma_E(s)\le C_Es\log s$ by Theorem~\ref{thm:main-group}(b). With $\epsilon=1/16$ the theorem gives
$Q\le C_Es\log s+2K_Gn/(s\ln2)+O(\log n)$ and $P\le2K_Gn/(s\ln2)+2$; take $s=\max(2K_G,\lceil\sqrt{n/\log n}\,\rceil)$.
\end{proof}

\subsection{Proof of Theorem~\ref{thm:main-channel}}\label{sec:main-channel-proof}

\emph{Closure and services.} By Theorem~\ref{thm:main-group}(b) the chunk $E$ of $\Gamma$ that consists of $1$ and the labels of the gadget of
$R_\Phi$ has $\log\sigma_E(r)\le C_Er\log r$. Proposition~\ref{prop:sofic-upper} therefore gives finite tracial channels within $K_G/r$ of $\Phi$ for
every $r$, so $\Phi\in\cK_{873}$; this also follows from~\cite{DouglasMghirbiB}, since $\Gamma$ is amenable (Lemma~\ref{lem:structure}) and hence
hyperlinear. Parts (c) and (a) of the proposition give the exact services and the zero-purity services in (a), since
$\log\sigma_E(R)\le C_ER\log R$ for $R\ge2$ and $\log(4K_Gn/B)\le(1+\log2K_G)\log(2n/B)$, because $\log(2n/B)\ge1$; for the zero-purity services use
$\epsilon=1/16$. Corollary~\ref{cor:streaming-lamps} gives the services in (b). The bounds on $r_*$ were proved in Section~\ref{sec:lamp-channel}.

\emph{Part (a).} Let a service have error $\epsilon\le1/16$ and $P+C\le B$, where $\log n\le B$. Theorem~\ref{thm:extraction} gives a round with
bath dimension $2^Q$, $\ddia(\Psi_t,\Phi)\le\epsilon/(1-\epsilon)\le1/15$ and $\nu_t=\ell_t+\Delta S_t\le t_n$. Since $\epsilon+h_2(\epsilon)<1$,
$\epsilon\log n\le\log n$ and $1\le\log n\le B$,
\[
t_n\le\frac{16}{15}\cdot\frac{B+1+\log n}{n}\le\frac{16B}{5n}\qquad(n\ge2).
\]
Let $C_R\ge1$ be the constant of Corollary~\ref{cor:amplification}(a) for $R_\Phi$ and the transport of Section~\ref{sec:lamp-channel}, and put
$y_n=5n/(16C_RB)$, so that $C_R\nu_ty_n\le1$ and $y_n\le n$. Put $c_0=5/(16C_Ry_0)$. If $B\le c_0n$, then $y_n\ge y_0$, and the corollary gives
\[
Q\ \ge\ S(\rho_t)\ \ge\ \frac{y_n}{32(\log y_n)^{2\kappa}}\ \ge\ \frac{5}{512\,C_R}\cdot\frac{n}{B(\log n)^{2\kappa}} .
\]
For the second statement take $B=\max(1,A_0)\log n$, which lies between $\log n$ and $c_0n$ for large $n$.

\emph{Part (b).} Proposition~\ref{prop:purity-budget}, with the flat probe of rank $r_*$, gives $P+C\le B:=3Q+6$. Theorem~\ref{thm:extraction} now
gives a round with $\nu_t\le t_n\le\frac{16}{15}(3Q+7+\log n)/n$. Put $y_n=15n/(16C_R(3Q+7+\log n))\le n$. If $y_n<y_0$, then
$3Q+7+\log n>15n/(16C_Ry_0)$, so $Q\ge c'n$ for large $n$. Otherwise Corollary~\ref{cor:amplification}(a) gives $Q\ge y_n/(32(\log n)^{2\kappa})$,
that is,
\[
Q\,(3Q+7+\log n)\ \ge\ \frac{15}{512\,C_R}\cdot\frac{n}{(\log n)^{2\kappa}} .
\]
If $Q\le\log n+7$ the left side is at most $4(\log n+7)^2$, which is smaller than the right side for large $n$. So $Q>\log n+7$, the left side is
at most $4Q^2$, and $Q\ge c\sqrt n/(\log n)^\kappa$. \qed

\begin{remark}[Memory between two readings of the profile]\label{rem:sandwich}
For every gadget channel of a sofic group that satisfies the hypotheses of Theorem~\ref{thm:weighted}, the least memory of services with error
$1/16$ and purity $O(\log n)$ lies between two readings of the group at scale $n$: below, the entropy forced by as many cheaply commuting copies of
$S_3$ as Corollary~\ref{cor:amplification} fits at relator defect about $\sqrt{\log n/n}$; above, $\log\sigma_E(Cn)+O(\log n)$ by
Proposition~\ref{prop:sofic-upper}, already at zero purity and with exchange $(n-1)Q$. For the configuration lamps both readings are
$n^{1+o(1)}$. The upper reading is the sofic profile, not the F{\o}lner function: the latter is doubly exponential for $\Gamma$
(Proposition~\ref{prop:folner-lower}), so models built from F{\o}lner sets~\cite[Proposition~3.14]{Cornulier2013} would only give memory
exponential in $n$, while the hash models
behind Theorem~\ref{thm:main-group}(b) give $O(n\log n)$. For the grid host of Remark~\ref{rem:grid-channel}, whose proofs are only sketched, the two readings would be $n/(\log n)^{15.75}$ and
$O(n\log n)$.
\end{remark}

\begin{remark}[The rank rate is the diagonal of the trade-off]\label{rem:sqrt-limit}
At the rank rate a device with error at most $1/16$ can spend at most $3Q+6$ bits of purity, so Theorem~\ref{thm:main-channel}(a) with $B$ of order $Q$ forces $Q\gtrsim\sqrt n$,
and the streaming services attain this with purity of the same order. In general, for every channel with a flat probe, padding a rank-rate
service with error at most $1/16$ with $2Q+6$ maximally mixed qubits turns it into a service whose purity is at most its memory; so the least memory at purity at most
memory is at most three times the least memory at the rank rate, plus six. For $\Phi$ both are $n^{1/2+o(1)}$. The one-round method cannot give
more than $\sqrt n$ at the rank rate, since condition~\eqref{eq:weighted-condition} certifies at most about $1/\nu$ copies; Theorem~\ref{thm:streaming}
shows that for $\Phi$ nothing more is true.
\end{remark}

%======================================================================
\section{Proof of the weighted one-round theorem}\label{sec:weighted-proof}
%======================================================================

The proof has four steps. First, a round with small leakage and entropy production yields unitaries on the bath under which every relation
holds up to $O(\sqrt\nu)$, with errors weighted by the memory state (Lemma~\ref{lem:weighted-extraction}). Second, weighted errors add up over van
Kampen diagrams, at a price for each face and each vertex that is paid once (Lemma~\ref{lem:charged-area}). Third, each transported copy of $S_3$ is rounded to an
exact representation on the same space, so the state and its entropy do not change (Lemma~\ref{lem:rounding}). Fourth, measuring the rounded
three-cycles one after another on one half of a purification of the memory state produces about one bit per copy that the other half also knows
(Theorem~\ref{thm:sector}).

The section generalises the argument of~\cite[Appendix~B]{DouglasMghirbiB}, which proves the one-round theorem for the Houghton
gadget~\cite[Theorem~5.5]{DouglasMghirbiB}, to every gadget channel that satisfies the hypotheses of Theorem~\ref{thm:weighted}. The weighted norms
are those of~\cite[Appendix~B]{DouglasMghirbiB}; Lemmas~\ref{lem:weighted-extraction} and~\ref{lem:charged-area} extend Lemmas~B.4 and~B.2
there from the Houghton gadget to every such gadget; Lemma~\ref{lem:rounding} is Lemma~B.3 there; and Theorem~\ref{thm:sector} is a form of
Theorem~B.1 there, proved by sequential measurement.

\subsection{Weighted norms}

Throughout, $\|\cdot\|_F$ is the unnormalized Hilbert--Schmidt norm on $M_M$, $\rho$ is a state on $\mathbb C^M$ and $\xi=\sqrt\rho$, so
$\|\xi\|_F=1$. For operators $W$ on $\mathbb C^M$ put
\[
E_\rho(W)=\|(W-I)\xi\|_F,\qquad\chi_\rho(W)=\|W\xi-\xi W\|_F,
\]
the \emph{displacement} and the \emph{centrality defect} of $W$ at $\rho$. We drop the subscript $\rho$. For unitaries $u,v$ and any operator $F$:
\begin{enumerate}[label=(W\arabic*),itemsep=1pt]
\item $E(uv)\le E(u)+E(v)$ and $E(u^*)=E(u)$;
\item $\chi(uv)\le\chi(u)+\chi(v)$ and $\chi(u^*)=\chi(u)$;
\item $\|Fv\xi\|_F\le\|F\xi\|_F+\|F\|\,\chi(v)$;
\item $|E(uWu^*)-E(W)|\le2\chi(u)$ for every $W$ with $\|W\|\le1$;
\item $\|\xi F\|_F=\|F^*\xi\|_F$, and $\|\xi F\|_F=\|F\xi\|_F$ if $F$ is normal.
\end{enumerate}
Indeed, $(uv-I)\xi=u(v-I)\xi+(u-I)\xi$ and $(u^*-I)\xi=-u^*(u-I)\xi$; $uv\xi-\xi uv=u(v\xi-\xi v)+(u\xi-\xi u)v$ and
$u^*\xi-\xi u^*=-u^*(u\xi-\xi u)u^*$; $Fv\xi=F\xi v+F(v\xi-\xi v)$; $(uWu^*-I)\xi=u(W-I)u^*\xi$, to which (W3) applies with $\|W-I\|\le2$, in both
directions; and $\|F^*\xi\|_F^2=\Tr(\xi FF^*\xi)$, which equals $\Tr(\xi F^*F\xi)$ when $F$ is normal. A word $w$ in $S^{\pm1}$ is evaluated by an
assignment $\phi$ of unitaries to $S$ through $\phi(x^{-1})=\phi(x)^*$; by (W2), $\chi(\phi(w))\le|w|\max_{x\in S}\chi(\phi(x))$.

\subsection{Extraction}

\begin{lemma}[Weighted extraction]\label{lem:weighted-extraction}
Let $G$, $S$ and $R$ be as in Section~\ref{sec:channel-setting}, with $a,b,j\in S$, $j^{-1}ab\in R$ and $j\ne1$ in $G$, and let $m$ be as in
Theorem~\ref{thm:weighted}. Let $\Psi(X)=\mathrm{Tr}_M\,U(X\otimes\rho)U^*$
be a round with $u=\ddia(\Psi,\Phi_G)\le1/8$, leakage $\ell$ relative to $\Phi_G$ and entropy production $\Delta S$, and put $\nu=\ell+\Delta S$,
$\eta=\sqrt{d\nu}$ and $e=256\,m\,\eta$. There are unitaries $\phi(x)$, $x\in S$, on $\mathbb C^M$ such that
\begin{enumerate}[label=(\roman*),itemsep=1pt]
\item $E(\phi(r))\le e$ for every $r\in R$;
\item $\chi(\phi(x))\le e$ for every $x\in S$;
\item $E(\phi(a)\phi(b))\ge1-e$.
\end{enumerate}
\end{lemma}

\begin{proof}
\emph{Normalization.} Let $(\zeta_i)$ and $(\zeta'_i)$ be orthonormal eigenbases of $\rho$ and $T(\rho)$ with decreasing eigenvalues $\lambda_i$ and $\mu_i$.
Replacing $U$ by $(I\otimes W)U$, where the unitary $W$ maps $\zeta'_i$ to $\zeta_i$, changes neither $\Psi$ nor its leakage and conjugates $T(\rho)$
by $W$. So we may assume $T(\rho)=\sum_i\mu_i|\zeta_i\rangle\langle\zeta_i|$. Input and output bath are now the same space, and
$\xi=\sum_i\sqrt{\lambda_i}|\zeta_i\rangle\langle\zeta_i|$ acts on both; this does not assert $T(\rho)=\rho$.

\emph{Entropy production controls commutators.} Write $U_{ji}=(I\otimes\langle\zeta_j|)U(I\otimes|\zeta_i\rangle)\in M_d$ and $t_{ji}=\frac1d\|U_{ji}\|_F^2$.
Since the columns and the rows of $U$ are orthonormal, $t$ is doubly stochastic, and the diagonal of $T(\rho)$ gives $\mu_j=\sum_it_{ji}\lambda_i$. The
distributions $p_{ji}=t_{ji}\lambda_i$ and $q_{ji}=t_{ji}\mu_j$ satisfy
\[
D(p\|q)=\sum_{j,i}t_{ji}\lambda_i\log\frac{\lambda_i}{\mu_j}=-S(\rho)+S(T(\rho))=\Delta S .
\]
By Jensen's inequality $-2\ln\sum\sqrt{p_{ji}q_{ji}}\le D(p\|q)\ln2$, and $-2\ln x\ge2(1-x)$, so
$\sum_{j,i}t_{ji}(\sqrt{\lambda_i}-\sqrt{\mu_j})^2\le\Delta S\ln2$. The matrix $t$ is a convex combination of permutation matrices (Birkhoff), and both sequences
decrease, so the rearrangement inequality gives $\sum_{j,i}t_{ji}\sqrt{\lambda_i\mu_j}\le\sum_j\sqrt{\lambda_j\mu_j}$; expanding the squares and using
double stochasticity, $\sum_j(\sqrt{\lambda_j}-\sqrt{\mu_j})^2\le\sum_{j,i}t_{ji}(\sqrt{\lambda_i}-\sqrt{\mu_j})^2$. The bath block $(j,i)$ of $U(I\otimes\xi)-(I\otimes\xi)U$
is $(\sqrt{\lambda_i}-\sqrt{\lambda_j})U_{ji}$, and $(\sqrt{\lambda_i}-\sqrt{\lambda_j})^2\le2(\sqrt{\lambda_i}-\sqrt{\mu_j})^2+2(\sqrt{\mu_j}-\sqrt{\lambda_j})^2$.
Therefore, with $U=\sum_{o,o'}|o\rangle\langle o'|\otimes A_{oo'}$ in system blocks $A_{oo'}\in M_M$,
\begin{equation}\label{eq:block-commutators}
\sum_{o,o'}\|[A_{oo'},\xi]\|_F^2=\|[U,I\otimes\xi]\|_F^2\le4d\,\Delta S\ln2=:\epsilon_S^2 .
\end{equation}

\emph{Leakage controls residuals.} The Choi state $\mathcal J(\Psi)$ is the marginal on $\mathbb C^d\otimes\mathbb C^d$ of the unit vector
$d^{-1/2}\sum_{o,o'}|oo'\rangle\otimes A_{oo'}\xi$ in $\mathbb C^d\otimes\mathbb C^d\otimes\mathcal H$, where $\mathcal H$ is $M_M$ with the Hilbert--Schmidt
inner product: indeed $\langle A_{o_2o_2'}\xi,A_{o_1o_1'}\xi\rangle=\Tr(A_{o_1o_1'}\rho A_{o_2o_2'}^*)=\langle o_1|\Psi(|o_1'\rangle\langle o_2'|)|o_2\rangle$. The support
of $\mathcal J(\Phi_G)$ is spanned by the orthogonal vectors $|A_g\rangle\!\rangle$, of squared norms $\mu_g=\|A_g\|_F^2$. Projecting onto it, put
\[
X_g=\mu_g^{-1}\sum_{[w_a]=g}\overline{c_a}A_a,\qquad R_{oo'}=A_{oo'}-\sum_g(A_g)_{oo'}X_g,
\]
where $A_a$ is the block at slot $a$. Then $\sum_{o,o'}\|R_{oo'}\xi\|_F^2=d\ell=:\epsilon_\ell^2$. At a slot $A_a=c_aX_{[w_a]}+R_a$, and at every other position
$A_{oo'}=R_{oo'}$. We have $\epsilon_\ell\le\eta$, $\epsilon_S\le2\sqrt{\ln2}\,\eta$, and, by Cauchy--Schwarz, $\epsilon_\ell+\epsilon_S\le\sqrt{1+4\ln2}\,\eta<2\eta$ and
$2\epsilon_\ell^2+2\epsilon_S^2\le8\ln2\,\eta^2$.

\emph{Anchors.} Every symbol $w$, and the empty word, occupies a diagonal position $z=z(w)$ whose row and column contain no other slot; there the
coefficient is $1$, so $B_w:=A_{zz}=X_{[w]}+R_{zz}$. Column unitarity gives $\Tr\rho(I-B_w^*B_w)=\sum_{o\ne z}\|R_{oz}\xi\|_F^2\le\epsilon_\ell^2$. Row
unitarity gives $\Tr\rho(I-B_wB_w^*)=\sum_{o'\ne z}\|\xi A_{zo'}\|_F^2\le\sum_{o'\ne z}(\|R_{zo'}\xi\|_F+\|[A_{zo'},\xi]\|_F)^2\le2\epsilon_\ell^2+2\epsilon_S^2$.
Let $B_w=V_w|B_w|=|B_w^*|V_w$ with $V_w$ unitary. Since $(1-x)^2\le1-x^2$ on $[0,1]$,
\[
\|(V_w-B_w)\xi\|_F\le\epsilon_\ell,\qquad\|\xi(V_w-B_w)\|_F=\|\xi(I-|B_w^*|)\|_F\le\sqrt{8\ln2}\,\eta,
\]
and so $\chi(V_w)\le\|(V_w-B_w)\xi\|_F+\|[B_w,\xi]\|_F+\|\xi(B_w-V_w)\|_F\le\epsilon_\ell+\epsilon_S+\sqrt{8\ln2}\,\eta<4.3\eta$, since
$\epsilon_\ell+\epsilon_S\le\sqrt{1+4\ln2}\,\eta$.

\emph{Triangle slots.} In the block of a triangle $(x,y,z)$, compare each slot $a$ with the anchor of the symbol $w'(a)$ that carries the same
label: the empty word at the upper left, $y$ at the upper right, $x$ at the lower left and $z$ at the lower right. Put $F_a=A_a-c_aV_{w'(a)}$. Then
$F_a=R_a-c_aR_{z'z'}+c_a(B_{w'(a)}-V_{w'(a)})$ with $z'=z(w'(a))$, and since $a$ and $(z',z')$ are different positions, Cauchy--Schwarz gives
\begin{gather*}
\|F_a\xi\|_F\le\sqrt{3/2}\,\epsilon_\ell+2^{-1/2}\epsilon_\ell<2\eta,\qquad\|F_a\|<2,\\
\|\xi F_a\|_F\le\|F_a\xi\|_F+\|[A_a,\xi]\|_F+2^{-1/2}\chi(V_{w'(a)})<7\eta .
\end{gather*}
For two slots $a_1,a_2$ in the same row of the block, with $A_{a_i}=c_iV_i+F_i$,
\[
A_{a_1}^*A_{a_2}-\overline{c_1}c_2V_1^*V_2=A_{a_1}^*F_2+c_2F_1^*V_2,
\]
whose right-weighted norm is at most $\|F_2\xi\|_F+2^{-1/2}(\|\xi F_1\|_F+\|F_1\|\chi(V_2))<2\eta+2^{-1/2}\cdot17\eta<15\eta$ by (W3) and (W5). The two
columns of the block are orthogonal columns of $U$, so the products for the two rows of the block add up to $-C_1^*C_2$, where $C_1,C_2$ are the
parts of the two columns in the other rows. These are residuals, so $\|C_1^*C_2\xi\|_F\le\|C_2\xi\|_F\le\epsilon_\ell$. The coefficients are
$\overline{c_1}c_2=\frac12$ for the pair (upper left, upper right) and $-\frac12$ for the pair (lower left, lower right), so
\[
\tfrac12\|(V_\varnothing^*V_y-V_x^*V_z)\xi\|_F<15\eta+15\eta+\eta .
\]
Put $W_w=V_\varnothing^*V_w$ for every symbol $w$, and $W_\varnothing=I$. Since $W_xW_y-W_z=V_\varnothing^*V_x(V_\varnothing^*V_y-V_x^*V_z)$,
\begin{equation}\label{eq:triangle}
\|(W_xW_y-W_z)\xi\|_F\le64\eta\quad\text{for every triangle }(x,y,z),\qquad\chi(W_w)<10\eta .
\end{equation}

\emph{Relations.} Let $\phi(x)=W_x$ for $x\in S$. For a letter $x\in S$, the triangle $(x,x^{-1},\varnothing)$ gives
$\|(W_{x^{-1}}-\phi(x^{-1}))\xi\|_F=\|W_x^*(W_xW_{x^{-1}}-I)\xi\|_F\le64\eta$. Fix a word $r=\sigma_1\cdots\sigma_\ell$ of $R$ and let
$D_l=\|(\phi(\sigma_1\cdots\sigma_l)-W_{p_l})\xi\|_F$, with $p_l$ its prefix symbols, $p_1=\sigma_1$ and $W_{p_\ell}=I$. Then $D_1\le64\eta$, and
\[
\phi(\sigma_1\cdots\sigma_l)-W_{p_l}=\bigl(\phi(\sigma_1\cdots\sigma_{l-1})-W_{p_{l-1}}\bigr)\phi(\sigma_l)+W_{p_{l-1}}\bigl(\phi(\sigma_l)-W_{\sigma_l}\bigr)+\bigl(W_{p_{l-1}}W_{\sigma_l}-W_{p_l}\bigr).
\]
By (W3), \eqref{eq:triangle} and the triangle $(p_{l-1},\sigma_l,p_l)$, $D_l\le D_{l-1}+2\cdot10\eta+64\eta+64\eta$. So
$E(\phi(r))=D_\ell\le148\,\ell\,\eta\le e$, which is (i); and (ii) holds because $\chi(\phi(x))<10\eta\le e$.

\emph{Displacement.} Let $z=z(j)$ and let $0$ be the position of the empty word. The labels there are $[j]\ne1$ and $1$, so
$\langle z|\Phi_G(|z\rangle\langle0|)|0\rangle=\tau(\lambda([j]))=0$, while $\langle z|\Psi(|z\rangle\langle0|)|0\rangle=\Tr(B_j\rho B_\varnothing^*)$. Feed
$(|zz\rangle+|00\rangle)/\sqrt2$ into $\Psi\otimes\mathrm{id}$ and $\Phi_G\otimes\mathrm{id}$. The difference $\Delta$ of the outputs is Hermitian and traceless with
$\|\Delta\|_1\le2u$, so each of its entries has modulus at most $\|\Delta\|\le\|\Delta\|_1/2\le u$; the entry at $(zz,00)$ is half the difference
above, so $|\Tr(\rho B_\varnothing^*B_j)|\le2u$. By the anchor bounds, $|\Tr(\rho V_\varnothing^*V_j)-\Tr(\rho B_\varnothing^*B_j)|\le2\epsilon_\ell$, and so
\[
E(\phi(j))^2=\|(V_j-V_\varnothing)\xi\|_F^2=2-2\,\mathrm{Re}\,\Tr(\rho V_\varnothing^*V_j)\ge2-4u-4\epsilon_\ell\ge\tfrac32-4\eta .
\]
By (i) for the word $j^{-1}ab$, $\|(\phi(a)\phi(b)-\phi(j))\xi\|_F=E(\phi(j)^*\phi(a)\phi(b))\le e$, so $E(\phi(a)\phi(b))\ge\sqrt{3/2-4\eta}-e\ge1-e$ if
$\eta\le1/8$. If $\eta>1/8$ then $e>1$ and (iii) is trivial.
\end{proof}

\subsection{Charged area}

Relator errors add up over a van Kampen diagram, but in the weighted norms every conjugation costs twice the centrality defect of the
conjugator. Writing a null word naively as a product of conjugates of relations would charge each conjugating word once for every relation it
encloses. Peeling the diagram one boundary face at a time charges instead each face a bounded amount and each vertex once.

\begin{lemma}[Charged area]\label{lem:charged-area}
Let $S$ and $R$ be finite, with every word of $R$ of length at most $m$, and let $\phi$ assign unitaries to $S$ with $E(\phi(r))\le e$ for $r\in R$
and $\chi(\phi(x))\le e$ for $x\in S$. Every word $w$ of finite area satisfies
\[
E(\phi(w))\ \le\ \bigl(3m\Area_R(w)+|w|\bigr)\,e\ \le\ \bigl((4m+1)\Area_R(w)+2|w|\bigr)\,e .
\]
\end{lemma}

\begin{proof}
Since $\phi$ is a homomorphism on the free group, $E(\phi(w))$ depends only on the free reduction of $w$, which is not longer and has the same
area; so let $w$ be freely reduced. By van Kampen's lemma~\cite[Chapter~V]{LyndonSchupp1977} there is a van Kampen diagram for $w$ with at most
$\Area_R(w)$ faces: a finite, connected and simply connected planar $2$-complex with edges labelled by letters of $S$, whose boundary cycle, read
from a base vertex, is $w$, and in which the boundary cycle of every face, read from a suitable vertex in a suitable direction, is a word of $R$.
We show, by induction on the number of edges, that every such diagram $D$ with $V$ vertices and $A$ faces satisfies
\begin{equation}\label{eq:charge}
E(\phi(w_D))\le\bigl((2m-1)A+2(V-1)\bigr)e ,
\end{equation}
where $w_D$ is its boundary word. If $D$ is a single vertex, $w_D$ is empty. Otherwise let $\varepsilon$ be the first edge of the boundary cycle,
from the base vertex $v_0$ to $v_1$, with label $x^{\pm1}$.

If $\varepsilon$ lies on no face, the boundary cycle traverses it twice, and removing it leaves diagrams $D_0\ni v_0$ and $D_1\ni v_1$ with
$w_D=x^{\pm1}w_{D_1}x^{\mp1}w_{D_0}$, the words read from $v_1$ and $v_0$. By (W1) and (W4), $E(\phi(w_D))\le2e+E(\phi(w_{D_1}))+E(\phi(w_{D_0}))$, and
the vertices and faces of $D_0$ and $D_1$ add up to those of $D$.

Otherwise $\varepsilon$ lies on the boundary of a face $f$. Let $c$ be the boundary cycle of $f$ read from $v_0$ and beginning with $\varepsilon$,
so that $c=\varepsilon\pi$ with $\pi$ a path in $\partial f$ from $v_1$ to $v_0$. Removing $f$ and $\varepsilon$ merges $f$ with the outer
region and leaves a van Kampen diagram $D'$ with the same vertices and one face fewer, whose boundary word read from $v_0$ is
$w_{D'}=\pi^{-1}\varrho$, where $w_D=\varepsilon\varrho$. Hence
$w_D=c\,w_{D'}$ freely. The word $c$ is a cyclic permutation of some $r^{\pm1}$, $r\in R$, that is $c=y^{-1}r^{\pm1}y$ freely with $|y|\le m-1$, so
$E(\phi(c))\le e+2(m-1)e$ by (W4) and (W2), and $E(\phi(w_D))\le(2m-1)e+E(\phi(w_{D'}))$ by (W1).

In both cases~\eqref{eq:charge} follows from the induction hypothesis. Finally $V-1$ is at most the number of edges of $D$, since $D$ is
connected, and counting the two sides of every edge gives twice the number of edges as at most $mA+|w|$, the sum of the lengths of the face
boundaries and of the boundary cycle. So~\eqref{eq:charge} gives $E(\phi(w))\le((2m-1)A+mA+|w|)e$.
\end{proof}

\subsection{Exact \texorpdfstring{$S_3$}{S3} rounding in place}

An approximate copy of $S_3$ is rounded to an exact one by functional calculus on the same space. The state is untouched, so no entropy is
gained or lost.

\begin{lemma}[{Exact rounding~\cite[Lemma~B.3]{DouglasMghirbiB}}]\label{lem:rounding}
Let $X$ and $Y$ be unitaries on $\mathbb C^M$, and suppose that the displacements of $X^2$, $Y^2$ and $(XY)^3$ and the centrality defects of
$X$ and $Y$ are at most $\delta$. There are unitaries $x,c$ on $\mathbb C^M$ with $x^2=c^3=I$ and $xcx=c^{-1}$, such that
\[
\|(x-X)\xi\|_F\le\delta,\qquad\chi(x)\le3\delta,\qquad\|(c-XY)\xi\|_F\le29\delta,\qquad\chi(c)\le56\delta .
\]
\end{lemma}

\begin{proof}
Let $\mathrm{sgn}\,t=1$ for $t\ge0$ and $-1$ for $t<0$, and let $x=\mathrm{sgn}(\mathrm{Re}\,X)$ and $y=\mathrm{sgn}(\mathrm{Re}\,Y)$ by functional calculus;
they are self-adjoint unitaries. On the unit circle $|z-\mathrm{sgn}\,\mathrm{Re}\,z|\le|z^2-1|$: for $z=e^{i\vartheta}$ with $|\vartheta|\le\pi/2$ this reads
$2|\sin(\vartheta/2)|\le2|\sin\vartheta|$, which holds as $\cos(\vartheta/2)\ge\frac12$, and the case $\mathrm{Re}\,z<0$ follows by symmetry. Hence
$|x-X|^2\le|X^2-I|^2$ as functions of the normal operator $X$, and $\|(x-X)\xi\|_F\le E(X^2)\le\delta$; by (W5) also $\|\xi(x-X)\|_F\le\delta$, so
$\chi(x)\le\chi(X)+2\delta\le3\delta$. The same holds for $y$.

Put $Z=XY$ and $w=xy$, so that $xwx=yx=w^*$. Then $\chi(Z)\le2\delta$, $\chi(w)\le6\delta$, and, by (W3),
$\|(w-Z)\xi\|_F\le\|x(y-Y)\xi\|_F+\|(x-X)Y\xi\|_F\le\delta+3\delta=4\delta$. From $w^3-Z^3=w^2(w-Z)+w(w-Z)Z+(w-Z)Z^2$ and (W3),
\[
E(w^3)\le E(Z^3)+4\delta+(4\delta+2\cdot2\delta)+(4\delta+2\cdot4\delta)\le25\delta .
\]
Let $f$ send each point of the unit circle to a nearest cube root of unity, with the ties $e^{\pm i\pi/3}$ and $-1$ sent to $1$. Then
$f(\bar z)=\overline{f(z)}$, and $|z-f(z)|\le|z^3-1|$: for $z=-1$ both sides equal $2$, and otherwise $z=f(z)e^{i\vartheta}$ with
$|\vartheta|\le\pi/3$, and the inequality reads $2|\sin(\vartheta/2)|\le2|\sin(3\vartheta/2)|$. Let $c=f(w)$. Then $c^3=I$, and $xcx=f(xwx)=f(w^*)=c^*=c^{-1}$. As before
$\|(c-w)\xi\|_F\le E(w^3)\le25\delta$ and $\|\xi(c-w)\|_F\le25\delta$. Hence $\|(c-Z)\xi\|_F\le29\delta$ and
$\chi(c)\le\chi(w)+\|(c-w)\xi\|_F+\|\xi(c-w)\|_F\le56\delta$.
\end{proof}

\subsection{A weighted sector bound}

Rounded copies of $S_3$ that nearly commute in the weighted sense force entropy. We work in the Hilbert space $\mathcal H$ of operators on
$\mathbb C^M$ with the Hilbert--Schmidt inner product $\langle Y,Z\rangle=\Tr(Y^*Z)$, identified with $\mathbb C^M\otimes\mathbb C^M$. Left multiplications
act on the first factor, which we call the system, and right multiplications on the second, the reference; they commute. The unit vector $\xi$
purifies $\rho$ on the system, and its reference marginal has the same nonzero spectrum, so it has entropy $S(\rho)$.

\begin{lemma}[Projection products~{\cite[Lemma~4.1]{DouglasMghirbiA}}]\label{lem:projection-product}
Let $P_1,\dots,P_k$ be orthogonal projections on a Hilbert space. Every vector $y$ satisfies
$\|y\|^2-\mathrm{Re}\langle y,P_k\cdots P_1y\rangle\le2\sum_l\|(I-P_l)y\|^2$.
\end{lemma}

\begin{proof}
By homogeneity let $\|y\|=1$. Put $y_0=y$, $y_l=P_ly_{l-1}$ and $s_l=y_{l-1}-y_l$, which lies in the range of $I-P_l$. Then
$\sum_l\|s_l\|^2=1-\|y_k\|^2=:\beta$, and $\alpha:=1-\mathrm{Re}\langle y,y_k\rangle=\mathrm{Re}\sum_l\langle(I-P_l)y,s_l\rangle\le\sqrt{E\beta}$ with
$E=\sum_l\|(I-P_l)y\|^2$. Since $0\le\|y-y_k\|^2=2\alpha-\beta$, we get $\alpha\le\sqrt{2E\alpha}$, that is $\alpha\le2E$.
\end{proof}

Let $c_1,\dots,c_N$ be unitaries on $\mathbb C^M$ with $c_i^3=I$, and write $c_i=\sum_{k\in\Z/3}\omega^kP_i^k$ with $\omega=e^{2\pi i/3}$ and spectral
projections $P_i^k$. For every $Y\in\mathcal H$,
\begin{equation}\label{eq:spectral-identities}
\|(c_i-I)Y\|_F^2=3\sum_{k\ne0}\|P_i^kY\|_F^2,\qquad\|[c_i,Y]\|_F^2=3\sum_{k\ne l}\|P_i^kYP_i^l\|_F^2,
\end{equation}
because $|\omega^k-\omega^l|^2=3$ for $k\ne l$ and the pieces $P_i^kYP_i^l$ are orthogonal. For $\mathbf a=(a_1,\dots,a_{i-1})\in(\Z/3)^{i-1}$ put
$Q_{\mathbf a}=P_{i-1}^{a_{i-1}}\cdots P_1^{a_1}$. Summing over the coordinates one at a time gives
$\sum_{\mathbf a}Q_{\mathbf a}^*Q_{\mathbf a}=\sum_{\mathbf a}Q_{\mathbf a}Q_{\mathbf a}^*=I$.

\begin{lemma}[Transfer to the reference]\label{lem:transfer}
For every $Y\in\mathcal H$ and $1\le i\le N$,
\[
\sum_{\mathbf a}\|Q_{\mathbf a}Y-YQ_{\mathbf a}\|_F^2\ \le\ \frac43\sum_{k<i}\|[c_k,Y]\|_F^2 .
\]
\end{lemma}

\begin{proof}
The pinchings $\mathcal E_k(Y)=\sum_lP_k^lYP_k^l$ are orthogonal projections on $\mathcal H$, and by~\eqref{eq:spectral-identities}
$\|(I-\mathcal E_k)Y\|_F^2=\sum_{l\ne l'}\|P_k^lYP_k^{l'}\|_F^2=\frac13\|[c_k,Y]\|_F^2$. Both $\sum_{\mathbf a}\|Q_{\mathbf a}Y\|_F^2$ and
$\sum_{\mathbf a}\|YQ_{\mathbf a}\|_F^2$ equal $\|Y\|_F^2$, and summing over the coordinates one at a time,
$\sum_{\mathbf a}\langle YQ_{\mathbf a},Q_{\mathbf a}Y\rangle=\sum_{\mathbf a}\Tr(Q_{\mathbf a}^*Y^*Q_{\mathbf a}Y)=\langle\mathcal E_1\cdots\mathcal E_{i-1}Y,Y\rangle$.
So the left side equals $2\|Y\|_F^2-2\,\mathrm{Re}\langle Y,\mathcal E_1\cdots\mathcal E_{i-1}Y\rangle$, and Lemma~\ref{lem:projection-product} bounds it by
$4\sum_{k<i}\|(I-\mathcal E_k)Y\|_F^2$.
\end{proof}

A three-valued distribution that is nearly symmetric under $k\mapsto-k$ carries about one bit for each unit of weight off $0$. Let $h_2$ be the
binary entropy. With $t=(1+y)/2$,
\[
\ln2\,\bigl(1-h_2(t)\bigr)=\sum_{k\ge1}\frac{y^{2k}}{2k(2k-1)}\ \le\ \sum_{k\ge1}\binom{2k}k\frac{y^{2k}}{(2k-1)4^k}=1-\sqrt{1-y^2}=1-2\sqrt{t(1-t)},
\]
comparing coefficients: $4^k\le2k\binom{2k}k$ holds for $k=1$, and the right side grows by the factor $2(2k+1)/k\ge4$ from $k$ to $k+1$. For a
distribution $(r_0,r_1,r_2)$ with $s=r_1+r_2>0$ this gives, with $t=r_1/s$,
\begin{equation}\label{eq:binary-hellinger}
H(r_0,r_1,r_2)\ \ge\ s\,h_2(r_1/s)\ \ge\ s-\frac{(\sqrt{r_1}-\sqrt{r_2})^2}{\ln2},
\end{equation}
which also holds when $s=0$.

\begin{theorem}[Weighted sector bound]\label{thm:sector}
Let $c_i,x_i$, $1\le i\le N$, be unitaries on $\mathbb C^M$ with $c_i^3=x_i^2=I$ and $x_ic_ix_i=c_i^{-1}$, with no condition relating different
indices. For a state $\rho$ with $\xi=\sqrt\rho$ put $\Delta_i=E(c_i)$, $\theta_i=\chi(c_i)$, $\theta'_i=\chi(x_i)$, and
\[
T_i=\Bigl(\sum_{k<i}\theta_k^2\Bigr)^{1/2},\qquad K_i=\Bigl(\sum_{k<i}\|[c_i,c_k]\xi\|_F^2\Bigr)^{1/2},\qquad K'_i=\Bigl(\sum_{k<i}\|[x_i,c_k]\xi\|_F^2\Bigr)^{1/2},
\]
\[
w_i=\frac13\Bigl(\Delta_i-\frac4{\sqrt3}T_i\Bigr)_+^2,\qquad\beta_i=\theta'_i+\frac2{\sqrt3}(2T_i+K'_i),\qquad
\pi_i=\frac13\Bigl(\theta_i+\frac2{\sqrt3}(2T_i+K_i)\Bigr)^2 .
\]
If $\pi_i\le2/3$ for all $i$, then
\[
S(\rho)\ \ge\ \sum_{i=1}^N\Bigl[w_i-\frac{\beta_i^2}{2\ln2}-h_2(\pi_i)-\pi_i\Bigr].
\]
\end{theorem}

\begin{proof}
\emph{Sequential measurement.} Measure $c_1,\dots,c_N$ in turn on the system, with the projections $P_i^k$, and let $A_i\in\Z/3$ be the outcomes.
Then $\Pr(A_{<i}=\mathbf a,A_i=k)=\|P_i^kQ_{\mathbf a}\xi\|_F^2=:p(\mathbf a,k)$. The measurements do not touch the reference, whose entropy is
$S(\rho)$, so
\[
S(\rho)\ \ge\ I(A_1\cdots A_N:\mathrm{Ref})=\sum_{i=1}^NI(A_i:\mathrm{Ref}\mid A_{<i}),
\]
since the conditional entropy of the reference given the classical outcomes is nonnegative. Fix $i$.

\emph{The copy stays displaced.} By~\eqref{eq:spectral-identities}, $\Pr(A_i\ne0)=\frac13\sum_{\mathbf a}\|(c_i-I)Q_{\mathbf a}\xi\|_F^2$. By
Lemma~\ref{lem:transfer} with $Y=\xi$, the vectors $Q_{\mathbf a}\xi$ and $\xi Q_{\mathbf a}$ differ by at most $\frac2{\sqrt3}T_i$ in $\ell^2$ over $\mathbf a$,
and $\sum_{\mathbf a}\|(c_i-I)\xi Q_{\mathbf a}\|_F^2=\Delta_i^2$ because $\sum_{\mathbf a}Q_{\mathbf a}Q_{\mathbf a}^*=I$. Since $\|c_i-I\|\le2$, the triangle
inequality in $\ell^2$ gives $\Pr(A_i\ne0)\ge w_i$.

\emph{The sign is a fair coin.} Since $x_iP_i^kx_i=P_i^{-k}$ and right multiplication by $x_i$ is unitary,
$\sqrt{p(\mathbf a,-k)}=\|P_i^kx_iQ_{\mathbf a}\xi\|_F$ and $\sqrt{p(\mathbf a,k)}=\|P_i^kQ_{\mathbf a}\xi x_i\|_F$, so
$\sum_{\mathbf a,k}(\sqrt{p(\mathbf a,k)}-\sqrt{p(\mathbf a,-k)})^2\le\sum_{\mathbf a}\|x_iQ_{\mathbf a}\xi-Q_{\mathbf a}\xi x_i\|_F^2$. Write
\[
x_iQ_{\mathbf a}\xi-Q_{\mathbf a}\xi x_i=x_i(Q_{\mathbf a}\xi-\xi Q_{\mathbf a})+\bigl((x_i\xi)Q_{\mathbf a}-Q_{\mathbf a}(x_i\xi)\bigr)+Q_{\mathbf a}(x_i\xi-\xi x_i).
\]
In $\ell^2$ over $\mathbf a$ the first term is at most $\frac2{\sqrt3}T_i$ and the third is $\theta'_i$. For the second apply Lemma~\ref{lem:transfer} with
$Y=x_i\xi$: since $[c_k,x_i\xi]=[c_k,x_i]\xi+x_i[c_k,\xi]$, it is at most $\frac2{\sqrt3}(K'_i+T_i)$. So the sum is at most $\beta_i^2$, and since the
terms $k=1$ and $k=2$ contribute equally, \eqref{eq:binary-hellinger} gives
\[
H(A_i\mid A_{<i})\ \ge\ \Pr(A_i\ne0)-\frac1{\ln2}\sum_{\mathbf a}\bigl(\sqrt{p(\mathbf a,1)}-\sqrt{p(\mathbf a,2)}\bigr)^2\ \ge\ w_i-\frac{\beta_i^2}{2\ln2}.
\]

\emph{The reference knows the outcome.} Let the reference measure $c_i$ from the right, with the projections $Y\mapsto YP_i^l$, and call the outcome
$B_i$. Left and right multiplications commute, so $\Pr(A_{<i}=\mathbf a,A_i=k,B_i=l)=\|P_i^kQ_{\mathbf a}\xi P_i^l\|_F^2$, and
by~\eqref{eq:spectral-identities} $\Pr(A_i\ne B_i)=\frac13\sum_{\mathbf a}\|[c_i,Q_{\mathbf a}\xi]\|_F^2$. As above,
\[
[c_i,Q_{\mathbf a}\xi]=c_i(Q_{\mathbf a}\xi-\xi Q_{\mathbf a})+\bigl((c_i\xi)Q_{\mathbf a}-Q_{\mathbf a}(c_i\xi)\bigr)+Q_{\mathbf a}(c_i\xi-\xi c_i),
\]
and Lemma~\ref{lem:transfer} with $Y=\xi$ and $Y=c_i\xi$ gives $\Pr(A_i\ne B_i)\le\pi_i$. By data processing and Fano's inequality for three values,
\[
I(A_i:\mathrm{Ref}\mid A_{<i})\ \ge\ I(A_i:B_i\mid A_{<i})\ \ge\ H(A_i\mid A_{<i})-h_2(\pi_i)-\pi_i ,
\]
where we used that $h_2(\pi)+\pi$ increases on $[0,2/3]$. Summing over $i$ proves the theorem.
\end{proof}

\begin{corollary}\label{cor:sector}
In Theorem~\ref{thm:sector}, suppose $\Delta_i\ge1/2$ for all $i$, and that $\theta_i$, $\theta'_i$, $\|[c_i,c_k]\xi\|_F$ and $\|[x_i,c_k]\xi\|_F$ are at
most $b$ for all $i\ne k$, where $\sqrt N\,b\le10^{-3}$. Then $S(\rho)\ge N/16$.
\end{corollary}

\begin{proof}
Now $T_i,K_i,K'_i\le\sqrt N\,b\le10^{-3}$ and $\theta_i,\theta'_i\le10^{-3}$. So $w_i\ge\frac13(\frac12-\frac4{\sqrt3}10^{-3})^2>0.0825$, and $\beta_i$ and
$\sqrt{3\pi_i}$ are at most $10^{-3}+2\sqrt3\cdot10^{-3}<4.5\cdot10^{-3}$. Hence $\beta_i^2/(2\ln2)<1.5\cdot10^{-5}$ and $\pi_i<6.8\cdot10^{-6}$, and with
$h_2(\pi)\le\pi\log(e/\pi)$ we get $h_2(\pi_i)<1.3\cdot10^{-4}$. Each summand exceeds $0.0823>1/16$.
\end{proof}

\subsection{Completion}

\begin{proof}[Proof of Theorem~\ref{thm:weighted}]
Let $\phi$ be given by Lemma~\ref{lem:weighted-extraction}, with $e=256\,m\sqrt{d\nu}$, and put
\[
b_*=100\,m\,(A+L+1)\,e=25600\,m^2(A+L+1)\sqrt{d\nu},\qquad\delta_L=(2L+1)e .
\]
Condition~\eqref{eq:weighted-condition} gives $\sqrt N\,b_*\le25600\cdot2^{-25}<10^{-3}$.

\emph{Copies.} Let $\psi_i=\phi(u_i)$, so $\chi(\psi_i)\le Le$, and put $X_i=\psi_i\phi(a)\psi_i^*$, $Y_i=\psi_i\phi(b)\psi_i^*$ and $Z_i=X_iY_i$. Since
$a^2$, $b^2$ and $(ab)^3$ lie in $R$, (W4) and Lemma~\ref{lem:weighted-extraction} give $E(X_i^2),E(Y_i^2),E(Z_i^3)\le e+2Le=\delta_L$, and (W2) gives
$\chi(X_i),\chi(Y_i)\le\delta_L$. Lemma~\ref{lem:rounding} gives unitaries $x_i,c_i$ with $x_i^2=c_i^3=I$, $x_ic_ix_i=c_i^{-1}$ and
\[
\|(x_i-X_i)\xi\|_F\le\delta_L,\qquad\chi(x_i)\le3\delta_L,\qquad\|(c_i-Z_i)\xi\|_F\le29\delta_L,\qquad\chi(c_i)\le56\delta_L .
\]
By (W4) and Lemma~\ref{lem:weighted-extraction}(iii), $E(Z_i)\ge E(\phi(a)\phi(b))-2Le\ge1-\delta_L$, so $\Delta_i=E(c_i)\ge1-30\delta_L$.

\emph{Raw commutators.} Let $i\ne k$ and $x,y\in\{a,b\}$, and let $P=\psi_i\phi(x)\psi_i^*$ and $P'=\psi_k\phi(y)\psi_k^*$. The group commutator
$g=[u_ixu_i^{-1},u_kyu_k^{-1}]$ is a word of length at most $8L+4$ and area at most $A$, so Lemma~\ref{lem:charged-area} gives
$E(\phi(g))\le(3mA+8L+4)e$. Since $PP'-P'P=(\phi(g)-I)P'P$ and $\chi(P'P)\le2\delta_L$, (W3) gives
\[
\|(PP'-P'P)\xi\|_F\le\gamma:=(3mA+8L+4)e+4\delta_L=(3mA+16L+8)e .
\]
Expanding $[X_i,X_kY_k]=[X_i,X_k]Y_k+X_k[X_i,Y_k]$ and $[X_iY_i,Z_k]=X_i[Y_i,Z_k]+[X_i,Z_k]Y_i$ and using (W3),
$\|[X_i,Z_k]\xi\|_F,\|[Y_i,Z_k]\xi\|_F\le2\gamma+2\delta_L$ and $\|[Z_i,Z_k]\xi\|_F\le4\gamma+6\delta_L$.

\emph{Rounded commutators.} For unitaries $U,V,U_0,V_0$,
\[
[U,V]-[U_0,V_0]=U(V-V_0)+(U-U_0)V_0-V(U-U_0)-(V-V_0)U_0,
\]
so by (W3) the weighted norms of the two commutators differ by at most
$2\|(U-U_0)\xi\|_F+2\|(V-V_0)\xi\|_F+2\chi(U_0)+2\chi(V_0)$. With $(U,U_0,V,V_0)=(x_i,X_i,c_k,Z_k)$ and $(c_i,Z_i,c_k,Z_k)$ this gives
\begin{gather*}
\|[x_i,c_k]\xi\|_F\le2\gamma+68\delta_L=(6mA+168L+84)e,\\
\|[c_i,c_k]\xi\|_F\le4\gamma+130\delta_L=(12mA+324L+162)e .
\end{gather*}
Since $m\ge6$, these and $\chi(c_i)\le(112L+56)e$ and $\chi(x_i)\le(6L+3)e$ are at most $b_*$, and $\Delta_i\ge1-(60L+30)e\ge1-b_*>1/2$.
Corollary~\ref{cor:sector} gives $S(\rho)\ge N/16$.
\end{proof}

The constant $2^{25}$ is not optimized; the proof uses only $25600\cdot2^{-25}<10^{-3}$, and Corollary~\ref{cor:sector} has room to spare.
